% GENERATED FILE. Edit canonical lessons, then run npm run generate:books.
% canonical-source-hash: fnv1a64:382072260c6f65df
% canonical-lessons: TA-C326-varalaru, TA-C326-katanta-kalam, TA-C326-palankalam, TA-C326-kalakattam, TA-C326-tacaptam

\chapter{The Past}
\label{ch:ta-the-past}
\hlchaptermodality{326}

\begin{chapteropening}
\textbf{By the end of this chapter:} \emph{I can talk about history and periods of time.}

Complete the last lesson of chapter 326: I can talk about history and periods of time.
\end{chapteropening}

\section[\texorpdfstring{\ta{வரலாறு} (varalāṟu)}{வரலாறு (varalāṟu)}]{\ta{வரலாறு} (varalāṟu) — history}

\label{lesson:TA-C326-varalaru}



\begin{quote}
Before the new one: say the Tamil for an epic, then the Tamil for to compose.
\end{quote}

\subsection*{You'll want to know: \ta{வரலாறு}}
\textbf{\ta{வரலாறு}} — \emph{varalāṟu} — \textquotedblleft{}history\textquotedblright{}.

\begin{quote}
\ta{எனக்கு} \ta{வரலாறு} \ta{படிக்க} \ta{மிகவும்} \ta{பிடிக்கும்}.
\end{quote}

\begin{quote}
\emph{eṉakku varalāṟu paṭikka mikavum piṭikkum.}
\end{quote}

\begin{quote}
I really like studying history.
\end{quote}

The oblique is \ta{வரலாற்று}-: \ta{வரலாற்றுப்} \ta{புத்தகம்}, a history book.

\begin{cousinweb}
\textbf{\ta{வரலாறு}} (\emph{varalāṟu}) is native Dravidian: DEDR 5270 lists it under \ta{வா}, \textquotedblleft{}to come\textquotedblright{}, as \textquotedblleft{}the order of events, history\textquotedblright{}: \ta{வரல்}, coming, and \ta{ஆறு}, way.
\end{cousinweb}

\begin{grammarlens}[title={What you've built}]
One more word for telling what happened.
\end{grammarlens}

\subsection*{Your turn}
\begin{itemize}
\raggedright
  \item \emph{Say it:} \emph{varalāṟu}
  \item \emph{Say it:} \emph{varalāṟu}, once more
  \item \emph{Say it:} \emph{eṉakku varalāṟu paṭikka mikavum piṭikkum.} — I really like studying history.
  \item \emph{Recall:} say the Tamil for an epic, then the Tamil for to compose, then say \emph{varalāṟu} again
\end{itemize}

\begin{tcolorbox}[breakable,colback=teal!4,colframe=teal!35!black,title={Before you move on}]
What does \ta{வரலாறு} mean? (\textquotedblleft{}History\textquotedblright{}.) Say it once more.
\end{tcolorbox}
\section[\texorpdfstring{\ta{கடந்த} \ta{காலம்} (kaṭanta kālam)}{கடந்த காலம் (kaṭanta kālam)}]{\ta{கடந்த} \ta{காலம்} (kaṭanta kālam) — the past}

\label{lesson:TA-C326-katanta-kalam}



\begin{quote}
Before the new one: say the Tamil for to compose, then the Tamil for history.
\end{quote}

\subsection*{You'll want to know: \ta{கடந்த} \ta{காலம்}}
\textbf{\ta{கடந்த} \ta{காலம்}} — \emph{kaṭanta kālam} — \textquotedblleft{}the past\textquotedblright{}.

\begin{quote}
\ta{அவர்} \ta{கடந்த} \ta{காலத்தைப்} \ta{பற்றி} \ta{அதிகம்} \ta{பேசுவதில்லை}.
\end{quote}

\begin{quote}
\emph{avar kaṭanta kālattaip paṟṟi atikam pēcuvatillai.}
\end{quote}

\begin{quote}
He does not talk much about the past.
\end{quote}

\ta{கடந்த}, gone by, from \ta{கட}, to pass.

\begin{cousinweb}
\textbf{\ta{கடந்த} \ta{காலம்}} (\emph{kaṭanta kālam}) joins \ta{கடந்த}, \textquotedblleft{}gone by\textquotedblright{}, from native \ta{கட}, \textquotedblleft{}to pass\textquotedblright{} (DEDR 1109), with \ta{காலம்}, \textquotedblleft{}time\textquotedblright{}, borrowed from Sanskrit \emph{kāla} (CDIAL 3084).
\end{cousinweb}

\begin{grammarlens}[title={What you've built}]
One more word for telling what happened.
\end{grammarlens}

\subsection*{Your turn}
\begin{itemize}
\raggedright
  \item \emph{Say it:} \emph{kaṭanta kālam}
  \item \emph{Say it:} \emph{kaṭanta kālam}, once more
  \item \emph{Say it:} \emph{avar kaṭanta kālattaip paṟṟi atikam pēcuvatillai.} — He does not talk much about the past.
  \item \emph{Recall:} say the Tamil for to compose, then the Tamil for history, then say \emph{kaṭanta kālam} again
\end{itemize}

\begin{tcolorbox}[breakable,colback=teal!4,colframe=teal!35!black,title={Before you move on}]
What does \ta{கடந்த} \ta{காலம்} mean? (\textquotedblleft{}The past\textquotedblright{}.) Say it once more.
\end{tcolorbox}
\section[\texorpdfstring{\ta{பழங்காலம்} (paḻaṅkālam)}{பழங்காலம் (paḻaṅkālam)}]{\ta{பழங்காலம்} (paḻaṅkālam) — ancient times}

\label{lesson:TA-C326-palankalam}



\begin{quote}
Before the new one: say the Tamil for history, then the Tamil for the past.
\end{quote}

\subsection*{You'll want to know: \ta{பழங்காலம்}}
\textbf{\ta{பழங்காலம்}} — \emph{paḻaṅkālam} — \textquotedblleft{}ancient times\textquotedblright{}.

\begin{quote}
\ta{பழங்காலத்தில்} \ta{மக்கள்} \ta{ஓலையில்} \ta{எழுதினார்கள்}.
\end{quote}

\begin{quote}
\emph{paḻaṅkālattil makkaḷ ōlaiyil eḻutiṉārkaḷ.}
\end{quote}

\begin{quote}
In ancient times people wrote on palm leaves.
\end{quote}

\ta{பழம்}, old + \ta{காலம்}; \ta{ம்} turns to \ta{ங்} before \ta{க}.

\begin{cousinweb}
\textbf{\ta{பழங்காலம்}} (\emph{paḻaṅkālam}) joins native \ta{பழ}, \textquotedblleft{}old\textquotedblright{} (DEDR 3999), with \ta{காலம்}, \textquotedblleft{}time\textquotedblright{}, borrowed from Sanskrit \emph{kāla} (CDIAL 3084).
\end{cousinweb}

\begin{grammarlens}[title={What you've built}]
One more word for telling what happened.
\end{grammarlens}

\subsection*{Your turn}
\begin{itemize}
\raggedright
  \item \emph{Say it:} \emph{paḻaṅkālam}
  \item \emph{Say it:} \emph{paḻaṅkālam}, once more
  \item \emph{Say it:} \emph{paḻaṅkālattil makkaḷ ōlaiyil eḻutiṉārkaḷ.} — In ancient times people wrote on palm leaves.
  \item \emph{Recall:} say the Tamil for history, then the Tamil for the past, then say \emph{paḻaṅkālam} again
\end{itemize}

\begin{tcolorbox}[breakable,colback=teal!4,colframe=teal!35!black,title={Before you move on}]
What does \ta{பழங்காலம்} mean? (\textquotedblleft{}Ancient times\textquotedblright{}.) Say it once more.
\end{tcolorbox}
\section[\texorpdfstring{\ta{காலகட்டம்} (kālakaṭṭam)}{காலகட்டம் (kālakaṭṭam)}]{\ta{காலகட்டம்} (kālakaṭṭam) — a period; an era}

\label{lesson:TA-C326-kalakattam}



\begin{quote}
Before the new one: say the Tamil for the past, then the Tamil for ancient times.
\end{quote}

\subsection*{You'll want to know: \ta{காலகட்டம்}}
\textbf{\ta{காலகட்டம்}} — \emph{kālakaṭṭam} — \textquotedblleft{}a period; an era\textquotedblright{}.

\begin{quote}
\ta{அந்தக்} \ta{காலகட்டத்தில்} \ta{கார்கள்} \ta{குறைவு}.
\end{quote}

\begin{quote}
\emph{antak kālakaṭṭattil kārkaḷ kuṟaivu.}
\end{quote}

\begin{quote}
In that period there were few cars.
\end{quote}

A -\ta{அம்} noun: \ta{காலகட்டத்தில்}, in that period.

\begin{cousinweb}
\textbf{\ta{காலகட்டம்}} (\emph{kālakaṭṭam}) is a Tamil compound built from \ta{காலம்}, \textquotedblleft{}time\textquotedblright{}, borrowed from Sanskrit \emph{kāla} (CDIAL 3084), and \ta{கட்டம்}, \textquotedblleft{}a stage, a block\textquotedblright{}.
\end{cousinweb}

\begin{grammarlens}[title={What you've built}]
One more word for telling what happened.
\end{grammarlens}

\subsection*{Your turn}
\begin{itemize}
\raggedright
  \item \emph{Say it:} \emph{kālakaṭṭam}
  \item \emph{Say it:} \emph{kālakaṭṭam}, once more
  \item \emph{Say it:} \emph{antak kālakaṭṭattil kārkaḷ kuṟaivu.} — In that period there were few cars.
  \item \emph{Recall:} say the Tamil for the past, then the Tamil for ancient times, then say \emph{kālakaṭṭam} again
\end{itemize}

\begin{tcolorbox}[breakable,colback=teal!4,colframe=teal!35!black,title={Before you move on}]
What does \ta{காலகட்டம்} mean? (\textquotedblleft{}A period; an era\textquotedblright{}.) Say it once more.
\end{tcolorbox}
\section[\texorpdfstring{\ta{தசாப்தம்} (tacāptam)}{தசாப்தம் (tacāptam)}]{\ta{தசாப்தம்} (tacāptam) — a decade}

\label{lesson:TA-C326-tacaptam}



\begin{quote}
Before the new one: say the Tamil for ancient times, then the Tamil for a period.
\end{quote}

\subsection*{You'll want to know: \ta{தசாப்தம்}}
\textbf{\ta{தசாப்தம்}} — \emph{tacāptam} — \textquotedblleft{}a decade\textquotedblright{}.

\begin{quote}
\ta{கடந்த} \ta{தசாப்தத்தில்} \ta{நகரம்} \ta{நிறைய} \ta{மாறியது}.
\end{quote}

\begin{quote}
\emph{kaṭanta tacāptattil nakaram niṟaiya māṟiyatu.}
\end{quote}

\begin{quote}
In the last decade the city changed a lot.
\end{quote}

Ten years; \ta{பத்தாண்டு} says the same in plain Tamil.

\begin{cousinweb}
\textbf{\ta{தசாப்தம்}} (\emph{tacāptam}) is borrowed from Sanskrit, \emph{daśa}, \textquotedblleft{}ten\textquotedblright{}, and \emph{abda}, \textquotedblleft{}a year\textquotedblright{} (Monier-Williams).
\end{cousinweb}

\begin{grammarlens}[title={What you've built}]
One more word for telling what happened.
\end{grammarlens}

\subsection*{Your turn}
\begin{itemize}
\raggedright
  \item \emph{Say it:} \emph{tacāptam}
  \item \emph{Say it:} \emph{tacāptam}, once more
  \item \emph{Say it:} \emph{kaṭanta tacāptattil nakaram niṟaiya māṟiyatu.} — In the last decade the city changed a lot.
  \item \emph{Recall:} say the Tamil for ancient times, then the Tamil for a period, then say \emph{tacāptam} again
\end{itemize}

\begin{tcolorbox}[breakable,colback=teal!4,colframe=teal!35!black,title={Before you move on}]
What does \ta{தசாப்தம்} mean? (\textquotedblleft{}A decade\textquotedblright{}.) Say it once more.
\end{tcolorbox}
